\documentclass[9pt, aspectratio=169]{beamer}

% --- Theme Selection ---
\usetheme{metropolis} 
% Alternative: \usetheme{Berlin}

% --- Color Customization (Purple & Blue) ---
\definecolor{DeepPurple}{HTML}{4B0082}
\definecolor{RoyalBlue}{HTML}{002366}
\definecolor{SoftPurple}{HTML}{9370DB}
\definecolor{CyanBlue}{HTML}{00BFFF}

\setbeamercolor{palette primary}{bg=RoyalBlue, fg=white}
\setbeamercolor{palette secondary}{bg=DeepPurple, fg=white}
\setbeamercolor{title}{fg=DeepPurple}
\setbeamercolor{progress bar}{fg=CyanBlue, bg=SoftPurple}

% --- Packages for Technical Content ---
\usepackage[utf8]{inputenc}
\usepackage{booktabs}
\usepackage{listings} % For code snippets
\usepackage{amsmath}
\usepackage{graphicx}
\usepackage{subfigure}

% --- Presentation Info ---
\title{The Identification of Exoplanets using Physics-Aware Machine Learning}
\subtitle{Mid-Semester Review}
\author{Esraaj Sarkar Gupta, Anirudh Puri, Suyash Prakash}
\date{\today}
\institute{Machine Learning and Pattern Recognition / Plaksha University}

\begin{document}

% Title Slide
\maketitle

% Table of Contents
\begin{frame}{Overview}
    \tableofcontents
\end{frame}

% Section : Literature Review
\section{Literature Review}
\begin{frame}{Standard ML Approaches: Probabilistic Validation}
    % SPEAKER NOTE: We begin with Armstrong et al. (2021), who applied highly rigorous probabilistic machine learning to the Kepler dataset to separate confirmed planets from false positives.
    
    \textbf{Armstrong et al. (2021)} utilized a suite of probabilistic classifiers on phase-folded light curves and derived vetting metrics.
    
    \vspace{0.2cm}
    \begin{itemize}
        \item \textbf{Methodology:} Deployed a Gaussian Process Classifier (GPC) alongside Random Forest (RFC), Extra Trees, and Multilayer Perceptron models.
        \item \textbf{Quantitative Success:} The models achieved exceptional accuracy, with both the GPC and RFC reaching an Area Under Curve (AUC) metric of 0.999. The GPC attained a highly precise mean log-loss per sample of 0.54.
    \end{itemize}

    \vspace{0.2cm}
    \begin{block}{The Limitation: Inherited Bias}
        % SPEAKER NOTE: Despite the mathematical rigor of the Gaussian Process, the model is only as unbiased as its training data.
        \begin{itemize}
            \item \textbf{Pipeline Reliance:} Approximately 50\% of the known KOI planets used in their training set were originally validated by a single algorithm (VESPA). Consequently, the ML models risk perpetuating existing validation biases rather than independently evaluating the geometry.
            \item \textbf{Outlier Failure:} Standard unconstrained models fail to confidently classify rare or anomalous candidates that are not heavily represented in the specific training space.
        \end{itemize}
    \end{block}
\end{frame}

\begin{frame}{Standard ML Approaches}
    % SPEAKER NOTE: While Armstrong focused on probability, Bhamare et al. (2021) approached the dataset as a strict binary classification problem, which exposed a critical flaw in unconstrained ML.

    \textbf{Bhamare et al. (2021)} evaluated standard machine learning architectures to classify Kepler objects strictly as 'CONFIRMED' or 'FALSE POSITIVE'.

    \vspace{0.2cm}
    \begin{itemize}
        \item \textbf{Methodology:} Evaluated models including Support Vector Machines (SVM), Random Forests, and AdaBoost using 19 reduced features from the DR25 catalogue. 
        \item \textbf{Quantitative Success:} The models achieved near-perfect statistical scores. AdaBoost reached a stratified F-1 score of 98.03\%, and a Majority-Vote ensemble achieved 98.28\%.
    \end{itemize}

    \vspace{0.2cm}
    \begin{block}{The Limitation: Target Leakage Problem}
        % SPEAKER NOTE: A closer look at feature importance reveals these models achieved high scores by cheating.
        \begin{itemize}
            \item \textbf{Target Leakage:} The feature set included pre-calculated categorical diagnostic flags, specifically the Centroid Offset (\texttt{koi\_fpflag\_co}) and Not Transit-Like (\texttt{koi\_fpflag\_nt}) flags.
            \item \textbf{Algorithmic Bias:} Instead of evaluating continuous physical parameters (like depth or planetary radius), the algorithms assigned the highest feature importance to these categorical flags, effectively bypassing the astrophysical features.
        \end{itemize}
    \end{block}
\end{frame}

\begin{frame}{The Shift Toward Physics-Informed ML}
    % SPEAKER NOTE: To solve the black-box problem, recent literature is shifting toward Physics-Informed Machine Learning, focusing on feature engineering rather than just algorithmic complexity.
    
    \textbf{Ford (2025)} reviewed the application of machine learning for mitigating stellar variability in Radial Velocity (RV) exoplanet surveys.

    \vspace{0.2cm}
    \begin{itemize}
        \item \textbf{The Challenge:} Modern RV surveys are limited by the host star's natural spectral variability (e.g., spots, pulsations), which obscures the signals of low-mass, Earth-like planets.
        \item \textbf{Algorithmic Comparison:} The study evaluated a spectrum of ML approaches, ranging from standard multilinear regression to complex Convolutional Neural Networks (CNNs).
    \end{itemize}

    \vspace{0.2cm}
    \begin{block}{The Value of Feature Engineering}
        % SPEAKER NOTE: The core finding of this paper perfectly aligns with our methodology: engineering the right physical constraints is more valuable than deploying the heaviest neural network.
        Ford concluded that relatively simple ML techniques, when paired with \textbf{well-engineered, physics-informed features}, match the performance of highly complex deep learning models. Crucially, they preserve the \textit{interpretability and explainability} required to establish the credibility of Earth-like planetary detections.
    \end{block}
\end{frame}

% Section: Dataset and Preprocessing 
\section{Dataset and Feature Preprocessing}

\begin{frame}{Nature of the Dataset \& Selection Criteria}
    \begin{itemize}
        \item \textbf{The Nature of the Dataset} 
        \begin{itemize}
            \item The dataset is an archival collection of Kepler Objects of Interest (KOIs).
            \item The transitions from raw Threshold Crossing Events (TCEs) which are potential planetary signals are examined to "Confirmed" or "False Positive" dispositions.
        \end{itemize}
        \vfill
        \item \textbf{The reason for selection of Kepler DR25 Selection}
        \begin{itemize}
            \item \textbf{Precision and Duration:} Kepler’s 4-year continuous periods of observations provide higher sensitivity to smaller planets compared to TESS’s 27-day observation periods.
            \item \textbf{Stellar Completeness:} The K2 and TESS missions are active which means the labels are prone to change whereas the Kepler DR25 is stable in nature.
        \end{itemize}
    \end{itemize}
\end{frame}

\begin{frame}{Data Collection}
    \begin{itemize}
        \item \textbf{Primary Source Files Used for Analysis}
        \begin{itemize}
            \item \texttt{q1\_q17\_dr25\_koi}: The primary dataset containing 49 observations, error bounds, and metadata parameters.
            \item \texttt{q1\_q17\_dr25\_sup\_koi}: The supplementary dataset containing 123 derived physics and refined planetary metrics.
        \end{itemize}
        \vfill
        \item \textbf{Data Integration}
        \begin{itemize}
            \item \textbf{Data Merging:} The primary dataset was \texttt{left-merged} on the supplementary one using the unique Kepler ID (\texttt{"kepid"}). This was done to ensure the primary observational data remained the basis to filter out stars without any detected signal and to ensure that each planet had a unique row.
            \item \textbf{Synchronization:} Observational data was integrated with derived features from the supplementary dataset.
        \end{itemize}
    \end{itemize}
\end{frame}

\begin{frame}{Data Collection Methodology}
    \begin{itemize}
        \item \textbf{Data Collection Methodology}
        \begin{itemize}
            \item \textbf{Kepler:} Differential time-series photometry via a 0.95m Schmidt telescope; monitored 150k stars in a fixed field.
            \item \textbf{K2:} Ecliptic plane surveying using solar radiation pressure for 3-axis stability following mechanical reaction wheel failure of Kepler.
            \item \textbf{TESS:} All-sky scanning via four wide-field cameras; focuses on high frequency monitoring of stars 30-100x brighter than Kepler targets.
        \end{itemize}
    \end{itemize}
\end{frame}

\begin{frame}{Features and Pre-processing Breakdown}
    \begin{itemize}
        \item \textbf{Data Scaling}
        \begin{itemize}
            \item \textbf{Dimensionality Reduction:} Out of the total 141 features that the dataset offers, 20 core parameters were deemed to be physically relevant.
            \item \textbf{Input for Model:} Final reduction was down to 10 physical features after removing text based or completely null columns, flags and unique IDs.
        \end{itemize}
        \vfill
        \item \textbf{Refinement}
        \begin{itemize}
            \item \textbf{Missing Values:} We removed rows with missing datapoints instead of filling in  (imputing) because of the extreme precision required for astrophysical calculations..
            \item \textbf{Feature Selection:} Excluded NASA's flags (\texttt{i.e."koi\_fpflag..."}) to force the model to rely entirely on physical features.
            \item \textbf{Class Balancing:} Entries without a definitive binary disposition (True Exoplanet/ False Positive) were separated entirely. 
            % SPEAKER NOTE: We intend to use the data entries without a definitive disposition as a final research problem. ~ 1500 entries.
        \end{itemize}
    \end{itemize}
\end{frame}

% Section : Physics Awareness
\section{Physics Awareness}
\begin{frame}{1. Transit Depth Consistency} % Analyzing the geometry of an eclipse
    %  SPEAKER NOTE: A light curve is a graph that plots the brightness (or flux) of an astronomical object over a period of time. In the context of exoplanets we expect to see a U shaped dip in the curve as the planet crosses in front of the star.

    % SPEAKER NOTE: The feature KOI depth represents the lowest point in the U curve.
    Typically the fraction of blocked stellar flux (the brightness of the star) is computed from a best-fit model produced by Mandel-Agol (2002).

    % SPEAKER NOTE: Mandel 2002 states that one may approximate a lightcurve by assuming that the local region of the star blocked by the planet has constant brightness.

    \begin{columns}[c] % The [c] vertically centers the content in both columns

        % --- LEFT COLUMN ---
        \begin{column}{0.48\textwidth}
            \begin{figure}
                \centering
                \includegraphics[width=\linewidth, height=0.5\textheight, keepaspectratio]{images/transitdepth.jpg}
                \caption{Theoretical Depth}
                \label{fig:transit_depth}
            \end{figure}
        \end{column}

        % --- RIGHT COLUMN ---
        \begin{column}{0.48\textwidth}
            \begin{figure}
                \centering
                \includegraphics[width=\linewidth, height=0.5\textheight, keepaspectratio]{images/kepler8b_lightcurve.png}
                \caption{Kepler-8b Data}
                \label{fig:kepler8b}
            \end{figure}
        \end{column}

    \end{columns}
    The primary characteristic of a transit lightcurve is stored in the feature \texttt{koi\_depth}, which is a measure of the lowest point of a transit lightcurve graph.
\end{frame}

\begin{frame}{1. Transit Depth Consistency (continued)}
    The assertion is made that since the telescope makes 2D measurements, the stellar flux must be directly related to the area of the star observed.

    \begin{equation}
        \text{Area}_\text{Star} = \pi R_*^2
    \end{equation}

    During transit, a local region of the star is blocked by the transiting body. We thus make the relation

    \begin{equation}
        \delta_\text{theoretical}   = \frac{\text{Area}_\text{Planet}}{\text{Area}_\text{Star}}
                                    = \frac{\pi R_p^2}{\pi R_*^2}.
    \end{equation}

    A new feature \texttt{delta\_delta} ($\Delta\delta$) where the observed drop in stellar flux $\delta_\text{obs}$ is anchored to $\delta_\text{theoretical}$.

    \begin{equation}
        \Delta\delta = \Bigg|\delta_\text{obs} - \bigg(\frac{R_p}{R_*}\bigg)^2 \Bigg|
    \end{equation}
    
\end{frame}
% References for Transit Depth Consistency
%   > https://arxiv.org/abs/astro-ph/0210099 (Mandel 2002)

\begin{frame}{2. Duration Consistency}
    Recall Kepler's Third Law
    % SPEAKER NOTES: The square of the planet's orbital period is directly propotionate to the cube of the planet's semi-major axis.
    \begin{equation}
        P^2 = \frac{4\pi^2}{G(M_* + M_p)}a^3
    \end{equation}

    \begin{figure}
        \centering
        \includegraphics[width=0.5\linewidth]{images/eccentricities.png}
        \caption{Eylen and Albrecht --- Results}
        \label{fig:placeholder}
    \end{figure}

    A 2015 publication from Eylen and Albrecht concluded that small planets in Kepler Multiplanet systems have low eccentricities $\approx 0.05$. Under this assertion, a circular theoretical anchor is produced.

    % SPEAKER NOTES: Solving Kepler's Third Law, we may find the value a -- the planet's semi-major axis. Since M_* and M_p are known from the dataset (we can compute M_star from data in the KOI dataset and we have an extra degeneracy from Newton's Law of Gravitation).

    % SPEAKER NOTES: There is well cited 2015 study that demonstrates that small planets in Kepler Multiplanet systems have low eccentricities, a Rayleigh distribution with a scale parameter of \sigma = 0.049 (basically a circle). Evident from the diagram, this is true for planets around the size of 1 earthmass.
    % SPEAKER NOTES: We thus use a circle (an ellipse of e = 0) to generate a theoretical anchor.

\end{frame}

\begin{frame}{2. Duration Consistency (Continued)}

    \begin{figure}
        \centering
        \includegraphics[width=0.6\linewidth]{images/ellipses.png}
        \caption{Comparing various eccentricities of ellipses}
        \label{fig:ellipses}
    \end{figure}
    
\end{frame}

\begin{frame}{2. Duration Consistency (Continued)}
    A theoretical anchor is built under the circular assumption ($e=0$).
    \begin{equation}
        T_\text{theoretical}    = \frac{\text{transit distance}}{\text{transit velocity}}
                                = \frac{2R_*}{\frac{2\pi a}{P}}
    \end{equation}
    
    The final feature produced quantifies the residue between the theoretically approximated transit time and the observed transit time, given by the intrinsic feature \texttt{koi\_duration}.

    \begin{equation}
        \Delta T = \Bigg| \texttt{koi\_duration} - \frac{PR_*}{\pi a}\Bigg|
    \end{equation}
\end{frame}
% CITATION: ECCENTRICITY FROM TRANSIT PHOTOMETRY: SMALL PLANETS IN KEPLER MULTI-PLANET SYSTEMS HAVE LOW ECCENTRICITIES (2015)

\begin{frame}{3. Stellar Density Consistency}
    The Photo-Eccentric Effect, introduced in 2012 by Dawson and Johnson allows the extraction of an exoplanet's orbital eccentricity from its photometric light-curve. This gives us a relationship between the stellar density for a perfectly circular orbit $\rho_\text{circ}$ and the true orbital ellipse with eccentricity e.
    
    % SPEAKER NOTES: We try not to go into the technicalities of this paper due to its mathematical complexity.
    \begin{equation}
        \frac{\rho_\text{circ}}{\rho_*} = \Bigg(\frac{1 + e\sin\omega}{\sqrt{1 - e^2}}\Bigg)^3
    \end{equation}

    % SPEAKER NOTES: By applying our baseline assumption of a circular orbit, the entire right side of the equation collapses to exactly 1. 
    Under our circular assumption ($e = 0$), the right side of the equation collapses to $1$. This mathematically requires that for a standard Kepler planet, the density implied by the transit must match the star's true spectroscopic density.
    
    \begin{equation}
        \frac{\rho_{\text{transit}}}{\rho_{\text{actual}}} \approx 1
    \end{equation}

    \vspace{0.3cm}
    % SPEAKER NOTES: Our final engineered feature is simply this ratio. If the Random Forest sees a ratio wildly deviating from 1, it flags the system as a geometric anomaly—typically a blended eclipsing binary False Positive.
    \begin{block}{The Final Feature}
        Our model evaluates the intrinsic ratio: \textbf{$\frac{\rho_{\text{transit}}}{\rho_{\text{actual}}}$}. Extreme deviations from $1.0$ flag a violation of the circular geometry, identifying false positives.
    \end{block}
\end{frame}

\section{Planned Work}

\begin{frame}{What's Next?}
    \begin{itemize}
        \item Build classical ML models on the data using (a) physical features and physically derived features in the dataset, and (b) the physics-informed engineered features derived from the dataset.
        \item Since the features to be considered can separate exoplanets from false-positives based on decision boundaries, it may be wise to consider the use of Random Forest models.
        \item We hope to identify more workable physically anchored features.
        \item We hope to test our model on more recent datasets such as TESS and K2 to test for robustness.
    \end{itemize}
\end{frame}

\begin{frame}[allowframebreaks]{References}
    \begin{thebibliography}{99}
        \setbeamertemplate{bibliography item}[text]
        
        \bibitem{Ford2025}
        Ford, E. B. (2025). 
        \newblock Enhancing Exoplanet Surveys via Physics-informed Machine Learning. 
        \newblock \emph{Proceedings of the International Astronomical Union}.

        \bibitem{Seager2003}
        Seager, S., \& Mallén-Ornelas, G. (2003).
        \newblock A Unique Solution of Planet and Star Parameters from an Extrasolar Planet Transit Light Curve.
        \newblock \emph{The Astrophysical Journal}, 585(2), 1038.

        \bibitem{VanEylen2015}
        Van Eylen, V., \& Albrecht, S. (2015).
        \newblock Eccentricity from Transit Photometry: Small Planets in Kepler Multi-planet Systems Have Low Eccentricities.
        \newblock \emph{The Astrophysical Journal}, 808(2), 126.

        \bibitem{Dawson2012}
        Dawson, R. I., \& Johnson, J. A. (2012).
        \newblock The Photo-eccentric Effect and Kepler Planets.
        \newblock \emph{The Astrophysical Journal}, 756(2), 122.

        \bibitem{Armstrong2021}
        Armstrong, D. J., et al. (2021).
        \newblock Exoplanet validation with machine learning: 50 new validated Kepler planets.
        \newblock \emph{Monthly Notices of the Royal Astronomical Society}, 504(4), 5327--5344.

        \bibitem{Plavalova2024}
        Plávalová, E., \& Rosaev, A. (2024).
        \newblock Classifications for Exoplanet and Exoplanetary Systems -- Could it be developed? I. Exoplanet classification.
        \newblock \emph{arXiv preprint arXiv:2409.09666}.

        \bibitem{Shallue2018}
        Shallue, C. J., \& Vanderburg, A. (2018).
        \newblock Identifying Exoplanets with Deep Learning: A Five-planet Resonant Chain around Kepler-80 and an Eighth Planet around Kepler-90.
        \newblock \emph{The Astronomical Journal}, 155(2), 94.

        \bibitem{Bhamare2021}
        Bhamare, A. R., Baral, A., \& Agarwal, S. (2021).
        \newblock Analysis of Kepler Objects of Interest using Machine Learning for Exoplanet Identification.
        \newblock \emph{2021 IEEE 12th International Conference on Computing Communication and Networking Technologies}.
    \end{thebibliography}
\end{frame}


\end{document}

https://www.cambridge.org/core/journals/proceedings-of-the-international-astronomical-union/article/abs/enhancing-exoplanet-surveys-via-physicsinformed-machine-learning/32B6B107092C4D24CF26BB2EF69831BC